\sgaStableTarget{sga6.e06.expose}\section*{\underline{Exposé VI}\\[3pt] \underline{\underline{THE @@M0@@ OF A PROJECTIVE BUNDLE:}}\\[3pt] \underline{\underline{COMPUTATIONS AND CONSEQUENCES}}}
% Exposé VI: reviewed local symmetric-algebra and Proj notation.
\renewcommand{\Sym}{\mathbf{S}}
\renewcommand{\Projz}{\underline{\operatorname{Proj}}_0}
\begin{center}
\underline{by Pierre BERTHELOT}
\end{center}

\sgaStableTarget{sga6.e06.sec.1}\section*{1. Computation of the @@M1@@ of a projective bundle: the case of locally free sheaves of finite type}

In this section, the notation @@M2@@ will denote the Grothendieck group of locally free sheaves of finite type on a scheme @@M3@@.

The aim of this section is to prove the following theorem:

\sgaStableTarget{sga6.e06.thm.1-1}\textbf{Theorem 1.1.} \uline{Let} @@M4@@ \uline{be a quasi-compact scheme}, @@M5@@ \uline{a} @@M6@@\uline{-Module locally free of rank} @@M7@@, @@M8@@ \uline{the associated projective bundle}; \uline{the structural morphism}: @@M9@@ \uline{defines a ring homomorphism}: @@M10@@ \uline{endowing} @@M11@@ \uline{with the structure of a} @@M12@@-\uline{algebra}. \uline{Let} @@M13@@ \uline{be the class of the sheaf} @@M14@@. \uline{Then} @@M15@@ \uline{is a free} @@M16@@-\uline{module with basis} @@M17@@.

\uline{Plan of the proof}:

@@M18@@) The powers of @@M19@@ generate @@M20@@ as a @@M21@@-module.

To prove this, we shall establish the following points:

\begin{enumerate}[label=\alph*)]
\item If @@M22@@ is a @@M23@@-Module locally free of finite type, there exists a graded @@M24@@-Module @@M25@@ of finite presentation, flat over @@M26@@, such that @@M27@@ (@@M28@@ being the symmetric algebra of @@M29@@; for the notation @@M30@@, see EGA II 8.12.1).

\item There exists a finite graded resolution of @@M31@@ by graded @@M32@@-Modules @@M33@@ locally free of finite type: @@M34@@.

\item The @@M35@@ admit a finite filtration whose quotients are of the form @@M36@@, where the @@M37@@ are locally free of finite type on @@M38@@.
\end{enumerate}

